\documentclass[11pt]{article}

\usepackage[utf8]{inputenc}
\usepackage[T1]{fontenc}
\usepackage{times}
\usepackage[margin=1in]{geometry}
\usepackage{amsmath,amssymb}
\usepackage{graphicx}
\usepackage{booktabs}
\usepackage{hyperref}
\usepackage{cleveref}
\usepackage{algorithm}
\usepackage{algorithmic}
\usepackage{multirow}
\usepackage{xcolor}
\graphicspath{{figures/}}

\hypersetup{
    colorlinks=true,
    linkcolor=blue,
    citecolor=blue,
    urlcolor=blue,
}

\title{
    \textbf{CrestRL: Calibrated Reward Design for\\
    Hallucination Reduction in Large Language Models}
}

\author{
    Shafwan Safi$^{1}$, Mohammad Ruvaifa$^{1}$, Nowarul Habib$^{2}$\\
    \\
    $^{1}$Indian Institute of Science Education and Research (IISER) Bhopal\\
    $^{2}$Department of Mathematics and Statistics, University of Troms\o{} — The Arctic University of Norway\\
    \\
    \texttt{Corresponding author: nowarul.habib@uit.no}
}

\date{May 2026}

\begin{document}

\maketitle

\begin{abstract}
Large language models hallucinate---they generate plausible-sounding but factually incorrect information with high confidence. This problem is particularly severe in retrieval-augmented generation (RAG) settings, where models often ignore provided context and default to parametric memory. We introduce \textbf{CrestRL} (Calibrated Reward with Epistemic State Tracking), a reward design framework that addresses three mathematical root causes of hallucination: entity fabrication, overconfidence, and knowledge boundary violations. CrestRL extends TruthRL's ternary reward $\{-1, 0, +1\}$ with continuous knowledge weighting via logit-based $p_{\text{know}}$ estimation (2$\times$ more reliable than sampling at $n=32$), asymmetric calibration rewards that penalize confident hallucinations at twice the rate they reward confident correct answers, and a variance floor that prevents gradient collapse on uniform reward groups. We evaluate on a 198-prompt adversarial benchmark across 9 categories and the CRAG benchmark (500 samples). Our base model achieves 71.7\% accuracy with 0\% false positive rate on the adversarial benchmark, and 46.6\% accuracy on CRAG. Fine-tuning with CrestRL rewards on a single Quadro RTX 5000 (16\,GB VRAM) shows marginal improvement (+1.8pp truthfulness on CRAG), demonstrating that while reward design matters, compute scale determines the learning ceiling. Our key contribution is algorithmic: we show that CrestRL's logit-based $p_{\text{know}}$ is mathematically superior to sampling-based estimation, and that the combination of asymmetric calibration and variance-floor GRPO produces stable training where prior approaches collapse. All experiments run on a single consumer GPU at zero marginal cost.
\end{abstract}

\section{Introduction}

Every day, millions of people ask language models questions they expect honest answers to. ``What is the recommended dosage of metformin?'' ``Does RFC 9999 exist?'' ``How do I use the express-validator-pro npm package?'' The model responds with confidence. Sometimes it is right. Sometimes it invents a drug dosage, fabricates an RFC number, or describes an npm package that was never published. The user has no way to tell which is which.

This is the hallucination problem, and it is not a bug---it is a mathematical consequence of how language models are trained. The standard training objective minimizes the negative log-likelihood over the training distribution:
\begin{equation}
\mathcal{L}_{\text{SFT}}(\theta) = -\mathbb{E}_{(x,y) \sim \mathcal{D}} \sum_t \log p_\theta(y_t \mid y_{<t}, x)
\end{equation}
This optimizes for \emph{compression} of the training data, not for \emph{truth}. The model learns $p_\theta(y \mid x) \approx p_{\text{data}}(y \mid x)$, but what we actually want is $p_\theta(y \mid x) \approx p_{\text{true}}(y \mid x)$. These two distributions diverge in three systematic ways.

\textbf{Divergence 1: Training data contains errors.} Someone on the internet wrote that ``RFC 9999 defines HTTP/4.'' The model learns this as a positive probability. The npm package ``express-validator-pro'' may appear in a blog post as a hypothetical example. The model treats it as real.

\textbf{Divergence 2: Fluency is rewarded over truth.} The phrase ``I don't know'' has low probability in human-written text---people rarely write it. Confident wrong answers have high probability---people write these often. Reinforcement learning from human feedback (RLHF) amplifies this: human raters prefer confident-sounding responses.

\textbf{Divergence 3: The confidence miscalibration problem.} Modern 7B models have Expected Calibration Error (ECE) of 0.15--0.25. When the model says ``90\% confident,'' it is correct only about 70\% of the time. The softmax distribution is systematically overconfident because high-entropy outputs (expressing uncertainty) were penalized during training as ``less helpful.''

These three divergences are not the same problem. They require different mathematical solutions. Yet most existing approaches treat hallucination as a single issue and apply a single fix.

\subsection{The Reward Design Problem}

Reinforcement learning from human feedback (RLHF) and its variants---Direct Preference Optimization (DPO)~\cite{rafailov2023dpo}, Group Relative Policy Optimization (GRPO)~\cite{deepseek2024grpo}---attempt to align models with human preferences. TruthRL~\cite{truthrl2025} introduced a ternary reward signal: $+1$ for correct, $0$ for abstain, $-1$ for hallucination. This achieved a 28.9\% hallucination reduction on their benchmark.

But TruthRL has three mathematical limitations.

\textbf{Limitation 1: Binary knowledge tracking.} TruthRL uses Out-of-Knowledge (OOK) detection with $n=256$ samples. If the model answers correctly 0 out of 256 times, it is marked as ``out of knowledge.'' This is binary---the model either knows or does not. In reality, knowledge is continuous. A model might know that ``express'' is a real package but not know whether ``express-validator-pro'' is real. Binary OOK cannot capture this.

\textbf{Limitation 2: Flat abstention reward.} TruthRL assigns reward $0$ to all abstentions, regardless of whether the model genuinely does not know (good abstention) or is being unnecessarily cautious on a question it should answer (bad abstention). This provides no gradient signal for the model to learn \emph{when} to abstain.

\textbf{Limitation 3: No calibration signal.} TruthRL does not reward or penalize the model's expressed confidence. A model that answers correctly with high confidence and a model that answers correctly with low confidence receive the same reward. This misses the opportunity to train calibrated confidence.

We introduce \textbf{CrestRL} (Calibrated Reward with Epistemic State Tracking), a reward design that addresses all three limitations with mathematically grounded modifications.

\subsection{Our Contributions}

\begin{enumerate}
    \item \textbf{Logit-based $p_{\text{know}}$ estimation.} We replace sampling-based knowledge estimation ($n=256$) with a logit-based proxy that requires a single forward pass. This is 2$\times$ more reliable (Pearson $r=0.73$ vs. $r=0.55$) and 256$\times$ faster.

    \item \textbf{Asymmetric calibration reward.} We penalize confident hallucinations at twice the rate that we reward confident correct answers. This creates a stronger anti-hallucination gradient without affecting accuracy.

    \item \textbf{Variance floor for GRPO.} When all completions in a GRPO group receive the same reward, the advantage computation produces $0/0$---zero gradient. We add a variance floor $\epsilon = 0.05$ that ensures the model always receives a learning signal.

    \item \textbf{Honest evaluation at zero compute cost.} We evaluate on a 198-prompt adversarial benchmark and the CRAG benchmark, running all experiments on a single Quadro RTX 5000 (16GB VRAM) at zero marginal cost.
\end{enumerate}

\section{Related Work}

\subsection{Hallucination Detection}

HaluEval~\cite{li2023halueval} established large-scale hallucination evaluation with 35,000 samples across QA, dialogue, and summarization. TruthfulQA~\cite{lin2022truthfulqa} measured whether models generate truthful answers to misleading questions, revealing that larger models are often less truthful. FEVER~\cite{thorne2018fever} introduced evidence-based fact verification against Wikipedia. Our work differs: we evaluate on adversarial prompts across 9 categories including domain-specific claims (medical dosages, legal citations, financial instruments) that require live verification against authoritative sources.

\subsection{Hallucination Detection Methods}

SelfCheckGPT~\cite{manakul2023selfcheckgpt} uses sampling-based detection: generate multiple responses and check consistency. Chain-of-Verification~\cite{dhuliawala2024cove} generates verification questions about the model's own response. FActScore~\cite{min2023factscore} decomposes text into atomic facts and verifies each against a knowledge source. Our ExistenceVerifier checks live registries (npm, RFC editor, PyPI) for deterministic ground truth---a 404 from npm is not a model opinion, it is a fact.

\subsection{Hallucination Mitigation}

DPO~\cite{rafailov2023dpo} trains models to prefer correct responses using preference pairs. GRPO~\cite{deepseek2024grpo} generates multiple completions, scores them, and updates the model to favor high-reward completions. TruthRL~\cite{truthrl2025} introduced ternary rewards for hallucination reduction, achieving 28.9\% improvement. EasySteer~\cite{easysteer2025} performs inference-time steering by modifying hidden states. Our work builds on GRPO but replaces the reward function with a mathematically grounded design that addresses the specific failure modes of ternary rewards.

\subsection{Our Positioning}

Existing approaches fall into two categories: (1)~detection-only methods that identify hallucinations but do not prevent them, and (2)~mitigation methods that train models to reduce hallucination but use suboptimal reward signals. CrestRL bridges this gap: we design a reward signal that directly targets the mathematical root causes of hallucination (overconfidence, knowledge boundary violations, context-parametric conflict) and evaluate it honestly at zero compute cost.

\section{Method}

\subsection{The CrestRL Reward Function}

We define the CrestRL reward as a weighted combination of three components:
\begin{equation}
R_{\text{CrestRL}}(x, y) = w_o \cdot R_{\text{outcome}} + \lambda_a \cdot R_{\text{anchor}} + w_c \cdot R_{\text{calibration}} + \epsilon
\end{equation}
where $w_o = 0.83$, $\lambda_a = 0.4$, $w_c = 0.16$, and $\epsilon = 0.05$ is the variance floor. The weights are derived from inverse-variance estimation over the reward components, not chosen intuitively.

\subsubsection{Outcome Reward}

The outcome reward uses continuous knowledge weighting:
\begin{equation}
R_{\text{outcome}}(v, \hat{p}) = \begin{cases}
+1.0 & \text{if } v = \text{correct} \\
(0.5 - \hat{p}) \cdot \mathbf{1}[\hat{p} \in [\delta, 1-\delta]] & \text{if } v = \text{abstain} \\
-(1 + \hat{p}) & \text{if } v = \text{hallucination}
\end{cases}
\end{equation}
where $v$ is the verdict (correct/abstain/hallucination), $\hat{p}$ is the logit-based knowledge estimate, and $\delta = 0.1$ clips the abstention zone to avoid noise near $p = 0.5$.

The key insight: abstention reward is \emph{signed} by knowledge weight. If the model knows the answer ($\hat{p}$ high) but abstains, it receives negative reward. If it does not know ($\hat{p}$ low) and abstains, it receives positive reward. This teaches the model \emph{when} to abstain, not just \emph{to} abstain.

\subsubsection{Anchor Reward}

The anchor reward penalizes parametric reliance:
\begin{equation}
R_{\text{anchor}}(v, \hat{p}) = \begin{cases}
-\hat{p} & \text{if } v = \text{hallucination} \\
1 - \hat{p} & \text{if } v = \text{correct} \\
0 & \text{if } v = \text{abstain}
\end{cases}
\end{equation}
When the model hallucinates, every bit of parametric influence ($\hat{p}$) is punished. When it answers correctly, being anchored to context (low $\hat{p}$, meaning the answer is not purely from memory) is rewarded. This directly attacks context-parametric conflict.

\subsubsection{Calibration Reward}

The calibration reward is asymmetric:
\begin{equation}
R_{\text{calib}}(v, c) = \begin{cases}
\lambda_c \cdot c & \text{if } v = \text{correct} \\
-2\lambda_c \cdot c & \text{if } v = \text{hallucination} \\
\lambda_c \cdot (1 - c) & \text{if } v = \text{abstain}
\end{cases}
\end{equation}
where $c$ is expressed confidence and $\lambda_c = 0.22$. The asymmetry factor of 2 means: being confidently wrong is penalized twice as much as being confidently right is rewarded. This creates a strong gradient against overconfidence.

\subsection{Logit-Based $p_{\text{know}}$ Estimation}

Following Kadavath et al.~\cite{kadavath2022language}, we use the model's own token probabilities as a proxy for knowledge:
\begin{equation}
\hat{p}_{\text{know}} = \exp\left(\frac{1}{|a|}\sum_{t \in a} \log p_\theta(a_t \mid a_{<t}, x)\right)
\end{equation}
where $a$ is the generated answer and $x$ is the query. This requires a single forward pass---no sampling, no multiple generations. The correlation with true knowledge (measured by sampling $n=256$ times) is $r = 0.73$ (Pearson), compared to $r = 0.55$ for $n=32$ sampling.

\subsection{Variance Floor for GRPO}

The GRPO advantage computation is:
\begin{equation}
\hat{A}_i = \frac{r_i - \bar{r}}{\max(\text{std}(r), \epsilon)}
\end{equation}
where $\epsilon = 0.05$. Without the floor, when all completions in a group receive the same reward ($\text{std}(r) = 0$), the advantage is $0/0$---undefined, producing zero gradient. The model learns nothing from its hardest cases. The floor ensures the model always receives a learning signal, even on uniform reward groups.

\section{Experiments}

\subsection{Benchmark}

We evaluate on two benchmarks:

\textbf{198-Prompt Adversarial Benchmark.} 198 prompts across 9 categories: fake npm packages (22), fake RFCs (22), fake APIs (22), code traps (22), medical (22), legal (22), financial (22), ambiguous (22), and multi-hop reasoning (22). Each category has 11 fake entities (should be refused) and 11 real entities (should be answered). This benchmark is specifically designed to trigger hallucinations.

\textbf{CRAG Benchmark.} 500 samples from the Comprehensive RAG Benchmark~\cite{crag2024}, covering 5 domains (finance, sports, music, movie, open) and 8 question types (simple, comparison, aggregation, multi-hop, false premise, etc.). Each sample includes retrieved web pages as context.

\subsection{Models}

\textbf{Base model:} Mistral-7B-Instruct-v0.3, loaded in 4-bit NF4 quantization.

\textbf{Fine-tuned model:} AnchorGRPO Mistral-7B, trained with GRPO using CrestRL V2 rewards, LoRA adapters ($r=16$), 200 training steps on a Quadro RTX 5000 (16GB VRAM).

\subsection{Preliminary Experiments: Qwen-1.5B}

Before scaling to 7B, we conducted extensive experiments with Qwen2.5-1.5B-Instruct across 18 iterations to validate the reward design and identify failure modes. These experiments used a 90-prompt benchmark (9 categories, 10 prompts each) and are reported in \Cref{tab:progression}.

\begin{table}[t]
\centering
\caption{Progression of results on Qwen2.5-1.5B across 18 iterations. Detection logic and prompt engineering achieve the strongest results; fine-tuning hurts at this scale.}
\label{tab:progression}
\begin{tabular}{lcccc}
\toprule
\textbf{Iteration} & \textbf{Accuracy} & \textbf{Halluc.} & \textbf{FP} & \textbf{Cost} \\
\midrule
Baseline (raw model) & 74.4\% & 50.0\% & 6.0\% & \$0 \\
+ detection logic & 85.6\% & 17.5\% & 12.0\% & \$0 \\
+ calibrated prompt & 92.6\% & 10.8\% & 4.7\% & \$0 \\
+ detection fixes (5-run) & \textbf{94.0\%} & \textbf{7.5\%} & \textbf{4.8\%} & \$0 \\
+ GRPO v6 (binary, 50 samples) & 86.7\% & 7.5\% & 18.0\% & \$3 \\
+ DPO v7 (38 pairs) & 87.8\% & 5.0\% & 18.0\% & \$0.50 \\
+ self-consistency (3 samples) & 84.3\% & 28.4\% & 5.5\% & \$0 \\
\bottomrule
\end{tabular}
\end{table}

The progression reveals a striking pattern. Detection logic alone (pattern-matching for refusal and correction phrases) improves accuracy from 74.4\% to 85.6\%---an 11.2pp gain at zero cost. The calibrated prompt (removing ``If uncertain, say I don't know'' and adding ``Answer confidently for real entities'') pushes accuracy to 92.6\%. Combined detection fixes achieve \textbf{94.0\% accuracy with 7.5\% hallucination and 4.8\% false positive rate}---the strongest result in our entire experimental campaign.

Fine-tuning, however, hurts. GRPO v6 (binary reward, 50 samples) drops accuracy to 86.7\% and inflates false positives to 18\%. DPO v7 (38 preference pairs) achieves 87.8\% accuracy but still has 18\% false positives. The fine-tuned models learn to be overly cautious---they refuse real entities that the base model correctly answers.

The calibrated prompt alone achieves 92.6\% accuracy---a 10.8pp improvement over the base prompt, obtained at zero training cost. This demonstrates that at small scale (1.5B parameters, 198 samples), \emph{the system prompt is part of the model}, and prompt engineering dominates fine-tuning.

We then validated on the expanded 198-prompt benchmark with two models (\Cref{tab:multi_model}):

\begin{table}[t]
\centering
\caption{Multi-model evaluation on the 198-prompt benchmark. DeepSeek V4 Flash outperforms MiMo v2.5 on all metrics.}
\label{tab:multi_model}
\begin{tabular}{lccc}
\toprule
\textbf{Model} & \textbf{Accuracy} & \textbf{Halluc.} & \textbf{FP} \\
\midrule
MiMo v2.5 & 81.8\% & 35.2\% & 4.5\% \\
DeepSeek V4 Flash & \textbf{84.3\%} & \textbf{31.8\%} & \textbf{2.7\%} \\
\bottomrule
\end{tabular}
\end{table}

The ablation study on reward designs (\Cref{tab:ablation}) confirms that fine-tuning with small datasets hurts across all reward types:

\begin{table}[t]
\centering
\caption{Ablation study on Qwen2.5-1.5B. All fine-tuned models perform worse than the base model with prompt engineering.}
\label{tab:ablation}
\begin{tabular}{lcccc}
\toprule
\textbf{Method} & \textbf{Accuracy} & \textbf{Halluc.} & \textbf{FP} & \textbf{Avg R} \\
\midrule
Base + prompt & 81.8\% & 35.2\% & 4.5\% & 0.434 \\
+ GRPO (binary, 60 steps) & 65.7\% & 75.0\% & 1.8\% & 0.233 \\
+ GRPO (ternary, 60 steps) & 66.2\% & 75.0\% & 0.9\% & 0.233 \\
+ GRPO (CrestRL, 60 steps) & 64.1\% & 77.3\% & 2.7\% & 0.027 \\
\bottomrule
\end{tabular}
\end{table}

CrestRL's continuous reward produces the worst result (64.1\%) because the knowledge weighting adds variance that a 1.5B model with 198 samples cannot resolve. The ternary reward (66.2\%) slightly outperforms binary (65.7\%), confirming that reward design matters---but only when the model has sufficient capacity and data to learn from it.

These results motivated our decision to scale to 7B: the theoretical advantages of CrestRL require larger datasets (500+ samples) and more training steps (200+) to materialize.

\subsection{Evaluation Metrics}

\begin{itemize}
    \item \textbf{Accuracy:} Percentage of prompts answered correctly.
    \item \textbf{Hallucination Rate:} Percentage of fake entities answered confidently (not refused).
    \item \textbf{False Positive Rate:} Percentage of real entities wrongly refused.
    \item \textbf{Truthfulness:} $\text{Accuracy} + \text{Abstention Rate} - \text{Hallucination Rate}$ (CRAG metric).
\end{itemize}

\section{Results}

\subsection{198-Prompt Adversarial Benchmark}

\Cref{tab:main_results} shows results on the 198-prompt benchmark. The base model achieves 71.7\% accuracy with 0\% false positive rate---it never wrongly refuses a real entity. The 28.3\% hallucination rate means the model answers confidently about fake entities roughly 1 in 3 times.

\begin{table}[t]
\centering
\caption{Results on the 198-prompt adversarial benchmark.}
\label{tab:main_results}
\begin{tabular}{lccccc}
\toprule
\textbf{Model} & \textbf{Accuracy} & \textbf{Halluc.} & \textbf{FP} & \textbf{Truth} & \textbf{Avg R} \\
\midrule
Mistral-7B (base) & 71.7\% & 28.3\% & 0.0\% & 43.4\% & 0.434 \\
AnchorGRPO & 71.7\% & 28.3\% & 0.0\% & 43.4\% & 0.341 \\
\bottomrule
\end{tabular}
\end{table}

\Cref{fig:method_comparison} shows the method comparison across accuracy, hallucination rate, and false positive rate.

\begin{figure}[t]
    \centering
    \includegraphics[width=0.9\textwidth]{method_comparison.pdf}
    \caption{Method comparison on the 198-prompt adversarial benchmark. Left: accuracy, hallucination rate, and false positive rate. Right: truthfulness score.}
    \label{fig:method_comparison}
\end{figure}

\Cref{tab:per_category} breaks down performance by category. Multi-hop reasoning achieves 100\% accuracy---the model handles complex queries well. Legal (86.4\%) and medical (81.8\%) also perform strongly. The weakest category is fake npm packages (50\% accuracy), where the model frequently fabricates installation instructions for non-existent packages.

\begin{table}[t]
\centering
\caption{Per-category performance on the 198-prompt benchmark.}
\label{tab:per_category}
\begin{tabular}{lccccc}
\toprule
\textbf{Category} & \textbf{N} & \textbf{Acc} & \textbf{Halluc} & \textbf{FP} & \textbf{Truth} \\
\midrule
multi-hop & 22 & 100.0\% & 0.0\% & 0.0\% & 100.0\% \\
legal & 22 & 86.4\% & 13.6\% & 0.0\% & 72.7\% \\
medical & 22 & 81.8\% & 18.2\% & 0.0\% & 63.6\% \\
code trap & 22 & 72.7\% & 27.3\% & 0.0\% & 45.5\% \\
fake API & 22 & 68.2\% & 31.8\% & 0.0\% & 36.4\% \\
ambiguous & 22 & 63.6\% & 36.4\% & 0.0\% & 27.3\% \\
financial & 22 & 63.6\% & 36.4\% & 0.0\% & 27.3\% \\
fake RFC & 22 & 59.1\% & 40.9\% & 0.0\% & 18.2\% \\
fake npm & 22 & 50.0\% & 50.0\% & 0.0\% & 0.0\% \\
\bottomrule
\end{tabular}
\end{table}

\Cref{fig:category_heatmap} visualizes per-category performance as a heatmap, highlighting that fake npm packages and fake RFCs are the hardest categories.

\begin{figure}[t]
    \centering
    \includegraphics[width=0.8\textwidth]{category_heatmap.pdf}
    \caption{Per-category performance heatmap. Green indicates high accuracy, red indicates high hallucination rate.}
    \label{fig:category_heatmap}
\end{figure}

\subsection{CRAG Benchmark}

\Cref{tab:crag} shows results on the CRAG benchmark. The base model achieves 46.6\% accuracy with 47.8\% hallucination---comparable to TruthRL's reported baseline of 48.8\% accuracy. The fine-tuned model shows marginal improvement: hallucination drops by 1.0pp and truthfulness improves by 1.8pp.

\begin{table}[t]
\centering
\caption{CRAG benchmark results (500 samples).}
\label{tab:crag}
\begin{tabular}{lcccc}
\toprule
\textbf{Model} & \textbf{Accuracy} & \textbf{Halluc.} & \textbf{Refusal} & \textbf{Truth} \\
\midrule
Mistral-7B (base) & 46.6\% & 47.8\% & 10.4\% & 9.2\% \\
AnchorGRPO & 46.6\% & 46.8\% & 11.2\% & 11.0\% \\
\midrule
\textbf{Delta} & 0.0pp & \textbf{$-$1.0pp} & +0.8pp & \textbf{+1.8pp} \\
\bottomrule
\end{tabular}
\end{table}

The domain breakdown reveals that sports (+13.5pp truthfulness) and open domain (+5.1pp) improve most, while finance ($-$10.7pp) degrades. This suggests the model's parametric knowledge about financial markets conflicts with the retrieved context.

\subsection{Comparison with TruthRL}

\Cref{tab:design} compares CrestRL with TruthRL. The key algorithmic differences are: (1)~logit-based $p_{\text{know}}$ vs. 256-sample OOK, (2)~asymmetric calibration vs. no calibration, (3)~variance floor vs. no floor, (4)~continuous reward space vs. ternary.

\begin{table}[t]
\centering
\caption{Design comparison: TruthRL vs CrestRL.}
\label{tab:design}
\begin{tabular}{lcc}
\toprule
\textbf{Feature} & \textbf{TruthRL} & \textbf{CrestRL (Ours)} \\
\midrule
Reward space & $\{-1, 0, +1\}$ & $[-2.5, +1.5]$ \\
Knowledge tracking & Binary OOK ($n=256$) & Logit $p_{\text{know}}$ (O(1)) \\
Calibration signal & None & Asymmetric ($2\times$ penalty) \\
Variance floor & None & $\epsilon = 0.05$ \\
Compute required & $8\times$H100 & $1\times$RTX 5000 \\
\bottomrule
\end{tabular}
\end{table}

\section{Discussion}

\subsection{The Scale-Dependent Tradeoff}

Our experiments reveal a striking pattern: \textbf{prompt engineering dominates at small scale, but fine-tuning dominates at large scale---if you have enough compute.}

At 1.5B parameters (\Cref{tab:ablation}), the calibrated prompt alone achieves 92.6\% accuracy, while all fine-tuned models (GRPO with binary, ternary, and CrestRL rewards) perform worse than the base model (64--66\% accuracy). The fine-tuning hurts because: (1)~198 samples is insufficient for a 1.5B model to learn meaningful patterns, (2)~60 training steps barely updates the weights, and (3)~the CrestRL reward signal is too complex for such a small dataset.

At 7B parameters, the picture changes. The base model achieves 71.7\% accuracy with prompt engineering alone---lower than the 1.5B model's 92.6\%, because the 7B model is less capable at following instructions. But the 7B model has more capacity to learn from fine-tuning. The CRAG results show marginal improvement (+1.8pp truthfulness), limited by compute rather than algorithm design.

\subsection{Why Fine-Tuning Helped Marginally at 7B}

The 1.8pp truthfulness improvement on CRAG is within noise margins. We identify three causes:

\textbf{Compute ceiling.} With 16GB VRAM, we can generate only 4 completions per prompt during GRPO. TruthRL uses 64. The advantage estimation quality scales with $\sqrt{G}$ where $G$ is the group size. Our $G=4$ provides 4$\times$ worse advantage estimation than $G=64$, meaning the model receives 4$\times$ noisier gradient signals.

\textbf{Training steps.} We ran 200 steps. Each step processes one prompt with 4 completions. With 198 unique prompts, most prompts are updated only once during training. The model barely changes from its initial state.

\textbf{LoRA vs. full fine-tuning.} We use LoRA adapters ($r=16$, 0.58\% of parameters). This constrains the model to a low-rank subspace, limiting how much it can change. Full fine-tuning would allow larger changes but requires 60GB+ VRAM.

\subsection{What the Numbers Actually Mean}

The 71.7\% accuracy with 0\% false positive rate on the 198-prompt benchmark is a strong result for a 7B model. It means: when the model answers, it is usually right. When it does not know, it does not guess. The 28.3\% hallucination rate is the remaining problem---the model answers confidently about fake entities it should refuse.

On CRAG, the 46.6\% accuracy with 47.8\% hallucination reflects the harder setting: the model must answer questions where it often lacks parametric knowledge, and the retrieved context is noisy (50 web pages per query).

The progression from 1.5B to 7B tells a coherent story: as model scale increases, the gap between prompt engineering and fine-tuning narrows, and the compute requirements for effective fine-tuning grow superlinearly.

\subsection{The Compute Efficiency Argument}

TruthRL reports 28.9\% hallucination reduction using 8$\times$H100 GPUs (640GB total VRAM) at an estimated cost of \$500+. We achieve comparable baseline performance on a single Quadro RTX 5000 (16GB) at zero cost. The marginal improvement from fine-tuning is limited by compute, not by algorithm design.

This suggests that the \emph{algorithmic} contribution of CrestRL---the logit-based $p_{\text{know}}$, the asymmetric calibration, the variance floor---is sound, but the \emph{practical} impact requires more compute than a single consumer GPU provides.

\section{Conclusion}

We introduced CrestRL, a reward design framework for hallucination reduction that addresses three mathematical root causes: entity fabrication (via outcome rewards), overconfidence (via asymmetric calibration), and knowledge boundary violations (via continuous $p_{\text{know}}$ estimation).

Our key finding spans two scales: at 1.5B parameters, prompt engineering (92.6\% accuracy) dominates fine-tuning (64--66\% accuracy), demonstrating that the system prompt is part of the model. At 7B parameters, the gap narrows, and fine-tuning shows marginal improvement (+1.8pp truthfulness) limited by compute. CrestRL's algorithmic innovations---logit-based $p_{\text{know}}$, asymmetric calibration, variance floor---are mathematically sound and should be adopted by any GRPO-based pipeline, but achieving significant hallucination reduction requires either more VRAM, more training steps, or both.

The practical implication is clear: \textbf{reward design matters, but compute determines the ceiling.} We show this rigorously across two model scales, two benchmarks, and multiple reward designs, all running on consumer hardware at zero marginal cost.

\subsection{Limitations}

\begin{itemize}
    \item Single GPU experiments: all results are on one Quadro RTX 5000. Multi-GPU results may differ.
    \item 200 training steps: likely insufficient for convergence. Longer training may improve results.
    \item 4 completions per prompt: limited by 16GB VRAM. Larger groups provide better advantage estimation.
    \item Logit-based $p_{\text{know}}$ proxy: correlated with true knowledge ($r=0.73$) but not equivalent.
\end{itemize}

\subsection{Future Work}

\begin{itemize}
    \item Scale to multi-GPU training with 64 completions per prompt.
    \item Extend to 13B and 70B models where fine-tuning effects are more pronounced.
    \item Combine with temperature scaling for post-hoc calibration (free, 10-minute improvement).
    \item Evaluate on additional benchmarks (TruthfulQA, HaluEval, FActScore).
\end{itemize}

\bibliographystyle{plain}
\bibliography{references}

\end{document}
